\chapter{Graph-Structured Lambda Theories}
\label{sec:gslt}
% ===================================================================

\begin{definition}[Graph-Structured Lambda Theory]
A \emph{graph-structured lambda theory} (GSLT) is a triple
$\GSLT = (\terms, \eqs, \rules)$ where:
\begin{enumerate}[label=(\roman*)]
  \item $\terms$ is a \emph{grammar of terms}, possibly including binders.
        We write $\terms(\GSLT)$ for the set of closed terms generated by
        the grammar.
  \item $\eqs$ is a set of \emph{equations on terms}: a smallest equivalence
        relation $\equiv$ on $\terms(\GSLT)$ that erases syntactic differences
        without semantic significance (e.g., $\alpha$-equivalence, commutativity
        and associativity of parallel composition).
  \item $\rules$ is a set of \emph{rewrite rules}: each rule $r \in \rules$ is
        a pair $(l_r, r_r)$ of terms (left-hand side and right-hand side) such
        that $l_r \rewrite_r r_r$ defines a reduction relation on
        $\terms(\GSLT)/{\equiv}$.
\end{enumerate}
The reduction relation generated by all rules is written $\rewrite$; its
reflexive-transitive closure is $\rewrite^{*}$.
\end{definition}

\begin{example}
The $\lambda$-calculus, the $\pi$-calculus, the Rho calculus, and the
ambient calculus are all naturally presented as GSLTs.  The grammar of terms
is the usual one; the equations include $\alpha$-equivalence and, in
concurrent calculi, structural congruence; the rules are the computational
reduction rules ($\beta$-reduction, communication, etc.).
\end{example}

\section{The Category of GSLTs}

\begin{definition}[Morphism of GSLTs]
Let $\GSLT = (\terms, \eqs, \rules)$ and $\GSLT' = (\terms', \eqs', \rules')$
be GSLTs.  A map $f : \terms(\GSLT) \to \terms(\GSLT')$ is a
\emph{morphism} $f : \GSLT \to \GSLT'$ if it preserves bisimilarity:
for all $u_1, u_2 \in \terms(\GSLT)$,
\[
  u_1 \bisim u_2 \;\Longrightarrow\; f(u_1) \bisim f(u_2).
\]
GSLTs and their morphisms form a category $\mathbf{GSLT}$.
\end{definition}

\begin{remark}
Bisimilarity here is the standard strong bisimilarity induced by the labeled
transition system (LTS) associated with $\GSLT$; see Section~\ref{sec:lts-hml}.
The category $\mathbf{GSLT}$ is the natural setting in which to state
universal properties of constructions such as the reversibility envelope.
\end{remark}

% ===================================================================
